\section{Results}
\label{sec:evaluation}

We ask three questions. Do agents that hold the money stay within it? Do they use the market that \sys{} gives them?
And how well do they split the money?

\subsection{Setup}
We ran two rounds. The first (2026-11-01) built the guards and measured them. The second (2026-11-02) added a third
model and an unreliable seller, measured the models again on a wider set of jobs, and tested the agents on more
projects.

\paragraph{Models and prices} The CEO is \texttt{gpt-6-luna} (\$0.10 / \$0.50 per million input / output tokens). The
first round's market has \texttt{gpt-6-luna} (the cheap model), \texttt{muse-spark-1.3} from Meta's Model API (\$1.25 /
\$4.25, the expensive model), and the \texttt{calc} and \texttt{csp} agents. The second adds \texttt{gpt-5-nano}
(\$0.05 / \$0.40) and \texttt{gridfast}. Prices are LiteLLM's list prices on the run dates. \texttt{gpt-5-nano} is the
cheapest per token, but it reasons for longer, so one attempt on it costs three times as much as one on
\texttt{gpt-6-luna} (\$0.0030 against \$0.0010). We tried other small models and dropped them (Gemma 4 31B, whose
hosted endpoint did not answer, and Muse Glimmer 30B, which costs more per token than \texttt{gpt-6-luna}). We left out
local models, since on the laptop we had (Apple M5, 24 GB) a 4B thinking model took six minutes on one arithmetic job.

\paragraph{Calibration} The first round sent 120 jobs (two per job type, level and difficulty) once to each option, for
\$0.18 on \texttt{gpt-6-luna} and \$3.18 on \texttt{muse-spark-1.3}. The outside agents are our own programs and cost
nothing real. The second round sent 240 jobs (four per group) three times to each option, for \$0.70 on
\texttt{gpt-6-luna}, \$2.17 on \texttt{gpt-5-nano} and \$20.32 on \texttt{muse-spark-1.3}. Table~\ref{tab:calibration}
shows where the second round's models differ, and where they cross over.

\begin{table}[!htb]
\centering
\caption{Share of attempts answered right in the second round's calibration, by level (1 to 5; 36 attempts each).}
\label{tab:calibration}
\footnotesize
\setlength{\tabcolsep}{4pt}
\begin{tabular}{@{}lllll@{}}
\toprule
\tablehead Job type & \texttt{gpt-5-nano} & \texttt{gpt-6-luna} & \texttt{muse-spark-1.3} & Outside agent \\
\midrule
Arithmetic & 100 100 100 81 \textbf{53} & 94 100 97 \textbf{64 25} & 100 100 100 92 92 & 100 100 100 100 \textbf{0} (\texttt{calc}) \\
Logic grids & 94 100 100 100 94 & 100 94 94 94 94 & 100 100 100 100 100 & 100 at every level (\texttt{csp}) \\
 & & & & 67 67 50 58 50 (\texttt{gridfast}) \\
Aggregation & 100 92 \textbf{44 25 0} & 92 86 64 53 \textbf{36} & 100 100 92 100 \textbf{47} & none \\
String pipelines & 100 94 92 64 \textbf{47} & 92 83 94 86 67 & 100 100 100 94 100 & none \\
\bottomrule
\end{tabular}
\end{table}

\paragraph{What these tables can show} Three of the seven validity checks fail on each table (\S\ref{sec:checks}), and
each failure limits what this section can claim. V5 fails on both (knowing the hidden difficulty is worth at most 0.009
of the points at $\rho = 0$), so these runs cannot show what adapting during the task is worth, for any strategy. We
compare the up-front and during-the-task strategies of Table~\ref{tab:twobytwo} on what they score and spend, not on how
well they adapt. V4 fails, so the expensive model solves nearly every job it is given (it drops to 60\% or below only on
level-5 aggregation). V3 fails on the second table, so its three models do not line up from weak to strong, since
\texttt{gpt-5-nano} beats \texttt{gpt-6-luna} on arithmetic and loses to it on aggregation and strings.

\paragraph{Runs} The first round drew ten projects at $\rho = 0.5$ and $\tau = 0.5$ (mean budget \$0.056), and ran
\sys{} v2 live twice on each of the first six (mean budget \$0.043). Its outside agents charged \$0.001 a message live,
against the $c_\mathrm{ref}$ of \$0.00023 that the replayed strategies paid. The second round draws twenty projects with
the same $\rho$ and $\tau$ (mean budget \$0.058) and runs every strategy three times on each. \sys{} v2 and the LLM split
use three seeds, and the replayed strategies use three of each job's recorded attempts. Its outside agents charge
$c_\mathrm{ref}$ (\$0.000215 a message, and a quarter of that for \texttt{gridfast}), live and in replay alike. The
solvers and fixed strategies are replayed on the results table. The LLM split makes its plan live and carries it out by
replay. Each is scored under two budget rules. Under the first hub's rule, a call only needs some money left and is
charged in full, even past the budget. Under v2's rule, a call is capped at what is left before it is made, as it is for
\sys{} v2. Replay cannot know what a capped call would have answered, so it counts it as no answer. In the first
round's live runs, all ten calls by hires that hit their cap returned nothing.

We threw away two attempts at the second round's live runs. The first sent nine processes at once, went over the
provider's limit of two million tokens a minute, and lost runs to refused calls. We made a rate-limited call wait and
retry (for up to 104 seconds), and no call in the runs we report failed for this reason. In the second attempt, the
project text put the market, the statistics and the answer format after the last job, inside that job's section, which
\texttt{--by-reference} hides from the CEO. The text now puts them before the jobs, for every version. In total we spent
\$33.21. The first round cost \$4.50. The second cost \$23.42 to calibrate (\$20.32 of it on \texttt{muse-spark-1.3},
and \$0.23 on pilots), \$4.30 for its live runs and \$0.99 for the two discarded attempts.

\subsection{What we found}
Overall, our experiments show the following results:

\textbf{The guards cap every hire.} The second round made 120 live runs of \sys{} v2 (two versions, twenty projects,
three seeds) with 551 hires. No hire spent more than 1.11 times its budget. Thirteen went over, by 9 to 11\%, each on a
single call, and all thirteen were aggregation jobs, whose inputs are long lists of records. A call's output is capped at
what the budget has left after the estimated cost of its input, so an overrun can only happen when the hub undercounts
the input's tokens by more than its 10\% margin. In the first round, none of 48 hires went over (Table~\ref{tab:ablation}).

\textbf{Each guard is needed.} In the first round we ran the same projects with \sys{} as first built (without the call
cap and the last turn), and with the call cap alone (Table~\ref{tab:ablation}). Without the cap, a budget was a limit
checked before a call, and the call's cost was charged after it. Hires given \$0.0005 spent up to 372 times that, the
overruns were charged to the CEO's budget, and three of six runs ended with no answer after spending 153\% of the
budget. With the cap alone, no hire went over, but the CEO reread every job, long documents included, on every turn,
and four of six runs ended without an answer (three because the CEO could not pay for its own answer, one at its
40-turn limit). With the last turn added, two of twelve runs were lost and the score rose to 0.566.

\begin{table}[!htb]
\centering
\caption{\sys{} with and without its guards, on the first round's six projects.}
\label{tab:ablation}
\small
\begin{tabular}{@{}lrrr@{}}
\toprule
\tablehead & As first built & Capped calls only & \sys{} v2 \\
\midrule
Runs & 6 & 6 & 12 \\
Mean score & 0.374 & 0.278 & \textbf{0.566} \\
Hires that spent past their budget & 9 of 26 & 0 of 25 & \textbf{0 of 48} \\
Largest overrun (spent / given) & 372$\times$ & 1.0$\times$ & \textbf{1.0$\times$} \\
Runs that ended without an answer & 3 of 6 & 4 of 6 & \textbf{2 of 12} \\
Spent / budget & 1.53 & 0.62 & \textbf{0.71} \\
Calls refused with a price quote & 0 & 19 & 25 \\
\bottomrule
\end{tabular}
\end{table}

\textbf{Keeping long inputs away from the CEO raises the agents' score by 0.136.} The second round ran \sys{} v2 as it
was, and with two changes. Long sections were hidden from the CEO (\texttt{request\_by\_\allowbreak reference}), and hires were
chosen by the success rates the CEO is shown (\texttt{success\_rates}). With both, the agents scored 0.599 against
0.463, a difference of $+0.136$ over the twenty projects (95\% interval $[+0.067, +0.216]$). The CEO reads the request's
outline (8.7 thousand characters on the first project) instead of the whole request (264 thousand), and its own calls
fell from 43\% of the budget to 32\%. Messages to outside agents now carry the puzzle the CEO names, so \texttt{csp}
answered 48 of its 49 messages right, where before, 24 of its 44 replies had been ``no grid puzzle in the text''. The
outside agents together answered 74 of 119 right, against 30 of 122. Runs where the CEO could not afford to read its
request once fell from 12 to 0, and lost runs from 17 to 11. We did not test the two changes separately. Every effect
above comes from hiding long sections, and what choosing by success rates adds on its own is not measured.

\textbf{With enough budget to run the CEO, the agents match the cheap model alone.} Table~\ref{tab:main} scores every
strategy under v2's budget rule, with the first rule's scores beside them. On the sixteen projects whose budget is
\$0.006 or more, the agents scored 0.737. That is about the same as \texttt{gpt-6-luna} alone (0.740; difference
$-0.003$, interval $[-0.038, +0.030]$) and the equal split (0.761; $-0.024$, $[-0.056, +0.007]$), and better than
random choice ($+0.067$, $[+0.005, +0.134]$). The agents also scored above the cascade, the LLM split and
\texttt{gpt-5-nano} alone, but by margins too small for sixteen projects to settle. They trail the knapsack ($-0.080$,
$[-0.129, -0.023]$), and they spent 3.9 times what the cheap model alone spent.

\textbf{The agents win on hard jobs that a specialist can do.} Two job types have an outside agent in the market
(\texttt{calc} for arithmetic, \texttt{csp} and \texttt{gridfast} for logic grids). On the sixteen projects, over the
hard jobs of those two types (26 jobs, three runs each), the agents answered 62 of 78 right (79\%), against 54 of 78
for \texttt{gpt-6-luna} alone (69\%). That is 24 against 18 on hard arithmetic, and 38 against 36 on hard grids. They
were ahead on six projects and behind on two, and the interval over projects only just reaches zero. They send grids to
\texttt{csp}, which answered all 41 of its messages on these projects right, and they give hard jobs a second hire more
often than easy ones (46 of the 81 jobs that went to more than one hire were hard, while hard jobs are 37\% of all
jobs). On the other jobs they do worse than the cheap model (226 of 306 right, 74\%, against 239, 78\%), and the two
cancel out to the matching overall score above.

\textbf{Below a minimum budget, the agents cannot work.} Four of the twenty projects have budgets of \$0.0009 to
\$0.0051, because their jobs are cheap to solve and $B = \tau\,C_\mathrm{full}$. On them, \texttt{gpt-6-luna} alone
scores 0.724, and the LLM split, whose one planning call comes out of the budget, scores 0.635. The agents score 0.046.
A CEO that hires, asks and fires pays for a turn at each step. All eleven runs the agents lost were on these four
projects, where the CEO came back empty after spending most or all of the budget on its own calls. On the other sixteen
they lost none of 48. Over all twenty projects, the agents' 0.599 is below the cheap model alone (0.736).

\begin{table}[!htb]
\centering
\caption{Strategies on the second round's twenty projects ($\rho = 0.5$, $\tau = 0.5$, mean budget \$0.058), sorted by
score under v2's budget rule, and on the sixteen projects whose budget is \$0.006 or more. Each runs three times per
project (sixty runs). ``Vs.\ Luna'' is spending divided by that of every job on \texttt{gpt-6-luna} (\$0.0079 a
project). The last column replays the same strategies under the first hub's rule, which charges a call in full even
past the budget. \sys{} v2's two versions run live, as it was and with long sections hidden (both during the task).}
\label{tab:main}
\footnotesize
\setlength{\tabcolsep}{3.6pt}
\begin{tabular}{@{}lrrrrrrr@{}}
\toprule
\tablehead Strategy & Score & Sixteen & Spent (\$) & Of budget & Vs.\ Luna & Own calls (\$) & First rule \\
\midrule
Oracle-full (knows outcomes) & 0.826 & 0.838 & 0.028 & 0.48 & 3.50$\times$ & 0 & 0.830 (0.49) \\
Knapsack (solver, up front) & 0.788 & 0.817 & 0.035 & 0.60 & 4.42$\times$ & 0 & 0.794 (0.68) \\
Every job on \texttt{gpt-6-luna} & 0.736 & 0.740 & 0.008 & 0.14 & 1.00$\times$ & 0 & 0.723 (0.15) \\
Equal split & 0.713 & 0.761 & 0.011 & 0.18 & 1.35$\times$ & 0 & 0.713 (0.18) \\
LLM split (LLM, up front) & 0.695 & 0.709 & 0.033 & 0.57 & 4.15$\times$ & 0.006 & 0.699 (0.69) \\
Cascade (solver, during) & 0.682 & 0.712 & 0.050 & 0.86 & 6.29$\times$ & 0 & 0.762 (1.58) \\
Every job on \texttt{gpt-5-nano} & 0.628 & 0.700 & 0.023 & 0.39 & 2.87$\times$ & 0 & 0.630 (0.41) \\
Random & 0.612 & 0.670 & 0.029 & 0.50 & 3.69$\times$ & 0 & 0.614 (1.05) \\
\textbf{\sys{} v2, long sections hidden} & \textbf{0.599} & \textbf{0.737} & \textbf{0.031} & \textbf{0.52} & \textbf{3.84$\times$} & 0.006 & live \\
Hybrid by level & 0.493 & 0.494 & 0.052 & 0.90 & 6.60$\times$ & 0 & 0.531 (2.14) \\
\sys{} v2 as it was & 0.463 & 0.579 & 0.036 & 0.61 & 4.47$\times$ & 0.025 & live \\
Every job on \texttt{muse-spark-1.3} & 0.342 & 0.427 & 0.054 & 0.93 & 6.82$\times$ & 0 & 0.447 (2.04) \\
\bottomrule
\end{tabular}
\par\smallskip{\footnotesize\raggedright First rule: score, with spent / budget in parentheses. ``Sixteen'': score on the
sixteen projects whose budget is \$0.006 or more.\par}
\end{table}

\textbf{The agents respond to prices.} A hire's call that its budget cannot pay for is refused before it reaches the
provider, and the refusal quotes the price. In the second round, 203 calls were refused this way in the version that
hides long sections (253 in the other). In 112 of them, a later hire in the same run was given more than the refused one
(156 in the other). Table~\ref{tab:refusal} follows one such exchange from a first-round log. The CEO called a level-4
logic grid easy and gave it to \texttt{gpt-6-luna} with \$0.001. The hire's first call reasoned up to the cap its budget
allowed and returned nothing. Its next call was refused, and so was the same call on the expensive model after the hub
moved it up. The CEO then gave five times as much to a new hire, and twelve times as much to a third, whose answer was
right and worth 4 of the project's 22 points.

\begin{table}[!htb]
\centering
\caption{A refusal with a price quote, and what followed, from the log of one first-round live run (project p2, second
run; job j4, a level-4 logic grid worth 4 of the project's 22 points).}
\label{tab:refusal}
\footnotesize
\begin{tabular}{@{}r p{1.7cm} p{11.6cm}@{}}
\toprule
\tablehead & Given & What happened \\
\midrule
1 & \$0.001 & The CEO declares the job easy, and the hub hires \texttt{grid4} on \texttt{gpt-6-luna}. \\
2 & & \texttt{grid4}'s first call reasons up to the cap its budget allows, spends \$0.0010 and returns nothing. \\
3 & & Its next call is refused before it reaches the provider. \\
4 & & The hub moves \texttt{grid4} up to \texttt{muse-spark-1.3}, and the refusal quotes the price: ``grid4's pool has
\$0.0000, and a call on meta/muse-spark-1.3 costs about \$0.0005 to read its messages and \$0.0011 for a short answer:
reserve more.'' \\
5 & \$0.005 & The CEO checks its own budget and hires \texttt{grid4b}, which answers that the clues are inconsistent
(wrong). \\
6 & \$0.012 & The CEO hires \texttt{solve4}, whose solution is right. \\
\bottomrule
\end{tabular}
\end{table}

\textbf{The agents choose the expensive model for the right jobs, but give it too little money.} On the sixteen
projects, 65\% of the agents' tries on \texttt{muse-spark-1.3} (68 of 104) were on hard jobs, and half came after a
cheaper model had tried the same job. Its 79 hires took 55\% of what the agents spent (\$0.99 of \$1.81) and returned 12
right answers, while its measured success rate on hard jobs is 87 to 100\%. The median hire was given \$0.005, but an
attempt on the expensive model costs a median of \$0.006 to \$0.025 depending on the job type (a mean of \$0.010 to
\$0.061). So 78 of the 79 hires had a call refused because their budget could not pay for it, after spending nearly all
of it first (a median of \$0.0046). The price quote comes after the money is gone. And at $\tau = 0.5$, the budget pays
for a full attempt on the expensive model on only a few jobs, so these hires buy calls that cannot finish.

\textbf{Splitting the money pays off once the budget has room, and the cheap model alone cannot use it.} The live runs
use $\tau = 0.5$. Replaying the results table at $\tau$ of 0.25, 1 and 2 (no new model calls; Figure~\ref{fig:sweep}
and Table~\ref{tab:sweep}) shows where a strategy can win. On the sixteen projects, \texttt{gpt-6-luna} alone scores
0.740 to 0.749 at every budget. It spends 31\% of the budget at $\tau = 0.25$ and 7\% at $\tau = 2$, and has nothing to
buy with the rest. The cascade adapts during the task, as the agents do. It tries the cheapest option expected to be
good enough, moves up when a grid's check fails, and asks a third option when the two cheapest disagree. It trails the
cheap model by 0.051 at $\tau = 0.25$ and leads it by 0.133 at $\tau = 2$ ($[+0.078, +0.192]$, ahead on 12 of 16
projects). The knapsack's lead grows from 0.029 to 0.162. The best single model also changes, since the expensive model
alone passes the cheap one at $\tau = 2$ (0.815 against 0.749). The knapsack, which plans up front, leads by more than
the cascade at every budget, so the sweep does not show that adapting beats planning (V5 fails, as above). What it
shows is that a budget with room rewards a strategy that uses it. The agents ran only at $\tau = 0.5$, where the
cascade still trails the cheap model and the agents' tries on the expensive model are underfunded. We did not run them
at a larger budget, so whether they follow the cascade there is not measured, but the sweep says that a larger budget is
where they should win.

\begin{figure}[!htb]
\centering
\begin{tikzpicture}
\begin{axis}[
  width=0.78\linewidth, height=6.4cm,
  xmode=log, log basis x=2, xmin=0.2, xmax=2.5,
  xtick={0.25,0.5,1,2}, xticklabels={0.25,0.5,1,2},
  ymin=0.3, ymax=0.95, ytick={0.3,0.4,0.5,0.6,0.7,0.8,0.9}, yticklabels={0.3,0.4,0.5,0.6,0.7,0.8,0.9},
  xlabel={Budget $\tau$ (times what the best answers cost)},
  ylabel={Score (share of points)},
  label style={font=\small\color{ink}},
  tick label style={font=\footnotesize\color{muted}},
  axis line style={draw=rulecolor,line width=0.6pt}, tick style={draw=none},
  xmajorgrids, ymajorgrids, major grid style={draw=rulecolor!70,line width=0.4pt},
  every axis plot/.append style={line width=1.5pt,line join=round,line cap=round,mark size=3pt},
  legend style={font=\footnotesize,draw=none,fill=none,at={(0.5,-0.24)},anchor=north,column sep=10pt,
    cells={anchor=west}},
  legend columns=3,
]
\addplot[color=vizblue,mark=*,mark options={draw=white,line width=0.8pt,fill=vizblue}]
  coordinates {(0.25,0.696) (0.5,0.712) (1,0.783) (2,0.882)};
\addlegendentry{Cascade (adapts during the task)}
\addplot[color=vizorange,mark=square*,mark options={draw=white,line width=0.8pt,fill=vizorange}]
  coordinates {(0.25,0.747) (0.5,0.740) (1,0.748) (2,0.749)};
\addlegendentry{Every job on \texttt{gpt-6-luna}}
\addplot[color=vizaqua,mark=triangle*,mark options={draw=white,line width=0.8pt,fill=vizaqua}]
  coordinates {(0.25,0.776) (0.5,0.817) (1,0.826) (2,0.910)};
\addlegendentry{Knapsack (plans up front)}
\addplot[color=vizyellow,mark=diamond*,mark options={draw=white,line width=0.8pt,fill=vizyellow}]
  coordinates {(0.25,0.350) (0.5,0.427) (1,0.562) (2,0.815)};
\addlegendentry{Every job on \texttt{muse-spark-1.3}}
\addplot[only marks,mark=o,mark size=5pt,draw=ink,line width=1.1pt] coordinates {(0.5,0.737)};
\addlegendentry{\sys{} v2 agents (live)}
\draw[ink,line width=0.5pt] (axis cs:0.5,0.737) ++(-0.06cm,-0.19cm) -- (axis cs:0.44,0.645)
  node[anchor=north,inner sep=1pt,font=\footnotesize,text=ink,align=center,fill=white]
  {\sys{} v2 agents, live: 0.737\\(matches \texttt{gpt-6-luna})};
\end{axis}
\end{tikzpicture}
\caption{Where splitting the money pays off. Scores on the sixteen projects as the budget grows eightfold, replayed
from the results table (three runs per project). The cheap model alone stays flat, while the cascade and the knapsack
turn the extra budget into points. The agents ran live at $\tau = 0.5$ only (ringed). Table~\ref{tab:sweep} gives the
numbers and intervals.}
\label{fig:sweep}
\end{figure}

\begin{table}[!htb]
\centering
\caption{Replayed strategies on the sixteen projects at four budgets ($B = \tau\,C_\mathrm{full}$, a mean of \$0.036,
\$0.072, \$0.144 and \$0.288 a project), three runs each under v2's budget rule. The agents ran live at $\tau = 0.5$
only. The last row is the cascade's lead over the cheap model alone, with its 95\% interval over projects.}
\label{tab:sweep}
\footnotesize
\setlength{\tabcolsep}{4pt}
\begin{tabular}{@{}lrrrr@{}}
\toprule
\tablehead Strategy & $\tau = 0.25$ & $\tau = 0.5$ & $\tau = 1$ & $\tau = 2$ \\
\midrule
Every job on \texttt{gpt-6-luna} & 0.747 & 0.740 & 0.748 & 0.749 \\
Every job on \texttt{gpt-5-nano} & 0.635 & 0.700 & 0.721 & 0.723 \\
Every job on \texttt{muse-spark-1.3} & 0.350 & 0.427 & 0.562 & 0.815 \\
Equal split & 0.748 & 0.761 & 0.791 & 0.808 \\
Cascade (solver, during) & 0.696 & 0.712 & 0.783 & \textbf{0.882} \\
Knapsack (solver, up front) & 0.776 & 0.817 & 0.826 & 0.910 \\
Oracle-full (knows outcomes) & 0.825 & 0.838 & 0.844 & 0.947 \\
\sys{} v2 agents (live) & & 0.737 & & \\
\midrule
Cascade minus \texttt{gpt-6-luna} & $-0.051$ & $-0.028$ & $+0.035$ & $\mathbf{+0.133}$ \\
 & $[-0.108, +0.004]$ & $[-0.088, +0.031]$ & $[-0.036, +0.100]$ & $[+0.078, +0.192]$ \\
\bottomrule
\end{tabular}
\end{table}

\textbf{What the agents still get wrong.} Besides the expensive model, two things separate them from the knapsack.
First, the CEO keeps some jobs for itself and answers them worse than a fresh call would. In the version that hides long
sections, it answered 169 jobs itself and got 85\% right, while \texttt{gpt-6-luna} given each job on its own gets 91\%.
The 223 jobs it handed to hires were right 67\% of the time, against 68\%. Second, the difficulty the CEO declares
barely matches the jobs' hidden difficulty (rank correlation 0.10 over 128 hires), and its own calls still take 32\% of
the budget. It rarely buys from the cheap, unreliable \texttt{gridfast} (seven messages, five of them right).

Holding the other strategies to v2's rule lowers the ones that overspent under the first rule. The cascade, the hybrid
by level and the expensive model alone, which spent 1.58 to 2.14 times their budget under the first rule, fell by
0.080, 0.038 and 0.105. Random choice and the knapsack fell by only 0.002 and 0.006. It raised the cheap model alone by
0.013, because refusing a call the budget cannot pay for saves the money for later jobs (on one small project, the first
rule let the cheap model's calls spend up to 2.9 times the budget, and it scored 0.4 to 0.6 less).
